\clearpage
\SourceLeaf{149}
\begin{center}
[IV.]

{\large [KRITISCHE BEMERKUNGEN ZUR METHODE DER KLEINSTEN QUADRATE.]}
\end{center}



\Needspace{6\baselineskip}
\begin{center}
\textsc{Gauss} an \textsc{Olbers}. Göttingen, 22. Februar 1819.
\end{center}

\noindent . . . . . . . . . . . Auch beschäftige ich mich gerade mit einer neuen Begründung der sogenannten Methode der kleinsten Quadrate. Meine erste Begründung setzt voraus, dass die Wahrscheinlichkeit des Beobachtungsfehlers $x$ durch $e^{-hhxx}$ dargestellt wird. Dann liefert jene Methode mit aller Strenge und in allen Fällen die wahrscheinlichsten Ergebnisse. Ist das Fehlergesetz unbekannt, so ist es \textls{unmöglich}, die wahrscheinlichsten Ergebnisse aus \textls{bereits gemachten} Beobachtungen anzugeben. \textsc{Laplace} hat die Sache von einer anderen Seite betrachtet und ein Prinzip gewählt, das ebenfalls zur\,Methode\,der\,kleinsten\ Quadrate führt\HistoricalNote{*)}{*)}{Und zwar unabhängig vom Fehlergesetz.}, wenn die Anzahl der Beobachtungen unendlich groß ist. Bei einer mäßigen Anzahl von Beobachtungen bleibt man dagegen völlig im Dunkeln, wenn das Fehlergesetz unbekannt ist. Selbst \textsc{Laplace} weiß für diesen Fall nichts Besseres zu sagen, als dass man auch hier die Methode\ der kleinsten\ Quadrate anwenden möge, weil sie ein bequemes Rechnen ermöglicht. Ich habe nun gefunden, dass man bei der Wahl eines etwas anderen Prinzips als desjenigen von \textsc{Laplace} (eines Prinzips, bei dem \textls{niemand} bestreiten kann, dass wir es mit mindestens ebenso viel Recht annehmen dürfen wie das von \textsc{Laplace}, und das meines Erachtens \textls{jeder Unvoreingenommene} für natürlicher erklären muss als das von \textsc{Laplace}), alle jene\space%
